% Ad Familiares 4.6 --- headnote
% ad-familiares-04-06, mid-April 45 BC, Atticus's Ficuleanum estate

\textit{Cicero to Servius Sulpicius Rufus, written at Atticus's
Ficuleanum estate about the middle of April 45 BC (Perseus:
\textit{in Attici Ficuleano medio m.~Apr.~a.~709 (45)}). Tullia
had died in mid-February; Cicero, unable to remain at Tusculum
where she had died and unwilling to be in Rome, had withdrawn to
Atticus's small property at Ficulea, north-east of the city,
where he spent much of the spring in retreat. This is his reply
to Servius's letter of consolation (\textit{Fam.}~4.5), arriving
some weeks after it had crossed from Athens. The pair is one of
the great correspondence-pairs in Latin literature, and the
bound-volume reader is expected to read the two letters together.}

\textit{The letter engages Servius's argument philosophically ---
Cicero acknowledges that he should bear his case as a man of
Servius's wisdom would think it ought to be borne --- but at the
same time admits, with a candour the project will not see him
match again on this question, that the philosophical apparatus he
has spent his life building does not get him through this moment.
The exempla Servius might have invoked --- Q. Fabius Maximus,
L. Aemilius Paullus, Sulpicius Galus, M. Porcius Cato Censorius
--- are exempla Cicero himself adduces in order to set them
aside: those men, having lost their sons, were still being upheld
by the standing they had in the commonwealth, and Cicero has no
such standing left. The home that once received him when he came
sorrowing from the Forum no longer holds his daughter; the
commonwealth that once received him when he came sorrowing from
home is no longer his to return to. The English of \S 2 should
sit at this load-bearing line: he is away from both, because
neither can console the grief he takes from the other.}

\textit{In \S 3 Cicero turns from grief to the immediate
political question --- how he and Servius are to live through a
time ``to be governed entirely by the will of one man.'' The
delicacy of the formulation is characteristic of the period: the
``man of prudence, of generous bearing, and not unfriendly to me,
and a particularly warm friend to you'' is Caesar, and the
question is not how to act but how to be permitted by his favour
to keep quiet. The letter ends as it began: with a request to
see Servius soon.}
